\section{Conclusions}
\label{sec:conclusion}

The controllability of an exothermic CSTR with a dynamic jacket, a rate-limited coolant valve and a fouling heat exchanger was examined by separating three limits: open-loop instability, actuator authority and cooling capacity. The main findings are the following.
\begin{enumerate}
    \item The benchmark reactor has a single, stable steady state for every admissible coolant flow at its nominal feed temperature. An open-loop unstable saddle at $\TSad$~K, with $\lambda_1=+\lamASad$~min$^{-1}$, exists only when the feed is cooled to $\TfSad$~K. Claims of a metastable unstable window must be checked against the steady-state branches of the actual parameter set.
    \item For the unstable reduced model, the bound $\tau_c<\tau_u-\theta$ is exactly the condition $K_ck>1$ on the static loop gain; it is necessary, and for proportional control of this reactor it was also sufficient over the tested range. With the SIMC integral time the stabilizing interval shrinks to $\tau_c<\tcMaxPI$~min, and the default $\tau_c=\theta$ is unstable. The reduced model carries the fundamental bound $\|S\|_\infty\ge e^{p\theta}=\msBound$, above $M_s=\msMax$; the third-order loop, whose only dead time is half a sample, admits $M_s=\msOpt$.
    \item Tight jacket regulation removes the instability: the concentration--temperature sub-block $\mathbf{A}_{12}$ is stable at the saddle, so the master loop of a cascade regulates a stable plant. The open-loop gain ratio overestimates the master gain by $\gainErrNom\%$ at the stable nominal point and has the wrong sign at the saddle; Eq.~\eqref{eq:k1_exact} gives the exact value.
    \item Against a $+10$~K coolant supply shock the default-SIMC cascade reduces the IAE by $\iaeRedSimcTjf\%$ relative to the SIMC single loop, but is worse than a PI optimized under $M_s\le\msMax$ over the same $\horizonShockWin$~min window ($\iaeTjfCascade$ versus $\iaeTjfSisoopt$~K\,min); against a feed temperature shock the gap is larger ($\iaeTfCascade$ versus $\iaeTfSisoopt$~K\,min). Static feedforward from $T_{jf}$ removes a further $\iaeRedFF\%$ of the cascade IAE and is the only element that beats the optimal single loop ($\iaeTjfCascadeff$ versus $\iaeTjfSisoopt$~K\,min); the cascade master loop was not optimized, so this comparison favours the single loop.
    \item Two fouling limits must be distinguished. The \emph{analytical linear boundary} is the capacity at which the frozen-$UA$ linearization of a fixed-gain loop loses the saddle: $\fstabCascade$--$\fstabSiso\%\,UA_0$ for the SIMC loops, far above $UA_{\mathrm{crit}}=\fcritSad\%$. It is a Hopf bifurcation (complex pair at $\pm j\,\wHopfSiso$--$\wHopfCascade$~rad/min) with $K_ck=\KcKCascade$--$\KcKSiso>1$ at the crossing, so it is not a loss of static gain. The \emph{practical loss of regulation} in fast transients occurs earlier: $\uaRunPctSadSiso\%\,UA_0$ for the SIMC single loop at $\kappa_{\mathrm{foul}}=\kappaFoul$~L/mol, and the departure approaches the linear boundary only for slow fouling (within $\gapSlowSiso$--$\gapSlowCascade$ points at $\kappa_{\mathrm{foul}}=\kappaSweepLo$~L/mol). The nominal point fails at the thermodynamic limit, $UA_{\mathrm{crit}}=\fcritNom\%\,UA_0$, and at $55\%\,UA$ the heat generated at the setpoint exceeds the maximum removal for any control law.
    \item The linear boundary belongs to the controller, not to the heat exchanger: the $M_s$-constrained PI holds the saddle down to $\fstabSisoopt\%\,UA_0$. Scheduling the slave gain as $K_{c2}\propto1/(T_j-T_{jf})$ moves the cascade boundary only from $\fstabCascade\%$ to $\fstabCascadegs\%$, because the cascade is lost through the master loop, whose static gain $K_{c1}K_1$ grows from $\KcOneKOneNom$ to $\KcOneKOneCascade$ as $\det\mathbf{A}_{12}\to0$.
\end{enumerate}
The consequence for operation is that, on an open-loop unstable branch, the margin against fouling is set by the sensitivity of the tuned loop and not by the heat exchanger. Scheduling or adapting the \emph{master} gain on the estimated $K_1(UA)$, optimizing the cascade under the same $M_s$ constraint, predictive control with an estimated heat-transfer capacity, and an experimental check of the fouling law are the natural next steps.
